% Part: computability
% Chapter: recursive-functions
% Section: composition

\documentclass[../../../include/open-logic-section]{subfiles}

\begin{document}

\olfileid{cmp}{rec}{com}
\olsection{ترکیب}

جے $f$ تے $g$ قدرتی عدداں دے دو اک-جگہ فنکشن
ہون، تاں اسی انہاں دی ترکیب بنا سکدے آں:
$h(x) = f(g(x))$۔ نواں فنکشن~$h(x)$، فنکشناں
$f$ تے~$g$ توں \emph{ترکیب (composition)}
راہیں متعین ہوندا اے۔ اسی ایس نوں اک توں ودھ
دلیلاں والے فنکشناں تک عام کرنا چاہندے آں۔

ایہ کرن دا اک طریقہ ایہ اے: فرض کرو $f$ اک
$k$-جگہ فنکشن اے، تے $g_0$، \dots،
$g_{k-1}$، کل $k$ فنکشن نیں، جیہڑے سارے
$n$-جگہ نیں۔ فیر اسی نواں $n$-جگہ فنکشن
$h$ ایویں متعین کر سکدے آں:
\[
h(x_0, \dots, x_{n-1}) =
f(g_0(x_0, \dots, x_{n-1}), \dots, g_{k-1}(x_0, \dots, x_{n-1}))
\]
جے $f$ تے سارے $g_i$ قابلِ حساب ہون، تاں
$h$ وی قابلِ حساب اے: $h(x_0,
\dots, x_{n-1})$ حساب کرن لئی پہلاں
$y_i = g_i(x_0, \dots, x_{n-1})$ دیاں
قیمتاں ہر $i = 0$، \dots،~$k-1$
واسطے حساب کرو۔ فیر ایہ قیمتاں $f$ وچ
پا کے $h(x_0, \dots, x_{n-1}) =
f(y_0, \dots, y_{k-1})$ حساب کرو۔

موجودہ فنکشناں توں نواں فنکشن حساب کرن دی ایہ
وضاحت شاید حد توں ودھ پابند لگے۔ اک گل ایہ
کہ کدے اسی فنکشن دیاں ساریاں دلیلاں نہیں
ورتدے، جیویں اسی $g(x, y, z) = \Succ(z)$
نوں~$\Add$ دی بنیادی بازگشتی تعریف وچ
ورتن لئی متعین کیتا سی۔ فرض کرو ساڈے کول ہیٹھ لکھے
فنکشن ورتن دی اجازت اے:
\[
\Proj{n}{i}(x_0, \dots, x_{n-1}) = x_i
\]
فنکشناں~$\Proj{k}{i}$ نوں
\emph{پروجیکشن (projection)} فنکشن آکھدے نیں:
$\Proj{n}{i}$ اک $n$-جگہ فنکشن اے۔ فیر
$g$ نوں ایویں متعین کیتا جا سکدا اے:
\[
g(x, y, z) = \Succ(\Proj{3}{2}(x, y, z)).
\]
اتھے $f$ دا کردار $1$-جگہ فنکشن $\Succ$
نبھاؤندا اے، ایس لئی $k=1$۔ تے ساڈے کول
اک $3$-جگہ فنکشن $\Proj{3}{2}$ اے،
جیہڑا $g_0$ دا کردار نبھاؤندا اے۔ نتیجہ
اک $3$-جگہ فنکشن اے، جیہڑا تیسری دلیل
دا جانشین واپس کردا اے۔

پروجیکشن فنکشن ساڈے نوں دلیلاں دی ترتیب
بدل کے یا اکّو دلیل نوں دوبارہ ورت کے وی
نویں فنکشن تعریف کرن دی سہولت دیندے نیں۔
مثال لئی فنکشن $h(x) = \Add(x, x)$ نوں
ایویں متعین کیتا جا سکدا اے:
\[
h(x_0) = \Add(\Proj{1}{0}(x_0),\Proj{1}{0}(x_0)).
\]
اتھے $k=2$، $n=1$، $f(y_0,y_1)$ دا
کردار $\Add$ نبھاؤندا اے، تے
$g_0(x_0)$ تے $g_1(x_0)$ دوہاں دا کردار
~$\Proj{1}{0}(x_0)$، یعنی اک-جگہ
پروجیکشن فنکشن (شناختی فنکشن وی آکھدے نیں)،
نبھاؤندا اے۔

جے $f(y_0, y_1)$ پہلے توں ساڈے کول فنکشن
ہووے، تاں فنکشن $h(x_0, x_1) = f(x_1, x_0)$
نوں ایویں متعین کر سکدے آں:
\[
h(x_0, x_1) = f(\Proj{2}{1}(x_0, x_1),\Proj{2}{0}(x_0, x_1)).
\]
اتھے $k=2$، $n = 2$، تے $g_0$ تے
$g_1$ دے کردار، بالترتیب، $\Proj{2}{1}$
تے~$\Proj{2}{0}$ نبھاؤندے نیں۔

تہانوں ایہ فکر وی ہو سکدی اے کہ $g_0$،
\dots،~$g_{k-1}$ ساریاں دی دلیل-گنتی
اکّو~$n$ ہونی ضروری اے۔ (یاد رکھو کہ
فنکشن دی \emph{دلیل-گنتی (arity)} اوہدیاں
دلیلاں دی تعداد اے؛ $n$-جگہ فنکشن دی
دلیل-گنتی~$n$ اے۔) پر پروجیکشن فنکشن
شامل کرن نال مطلوبہ لچک مل جاندی اے۔
مثال لئی فرض کرو $f$ تے~$g$، $3$-جگہ
فنکشن نیں، تے $h$~اوہ $2$-جگہ فنکشن اے
جیہڑا ایویں متعین اے:
\[
h(x,y) = f(x,g(x,x,y),y).
\]
فنکشن~$h$ دی تعریف پروجیکشن فنکشناں نال
ایویں دوبارہ لکھ سکدے آں:
\[
h(x,y) = f(\Proj{2}{0}(x,y), g(\Proj{2}{0}(x,y), \Proj{2}{0}(x,y),
\Proj{2}{1}(x,y)), \Proj{2}{1}(x,y)).
\]
فیر $h$، $f$ دی $\Proj{2}{0}$،
$l$ تے $\Proj{2}{1}$ نال ترکیب اے، جتھے
\[
l(x,y) = g(\Proj{2}{0}(x,y),\Proj{2}{0}(x,y),\Proj{2}{1}(x,y)),
\]
یعنی $l$، $g$ دی $\Proj{2}{0}$،
$\Proj{2}{0}$ تے~$\Proj{2}{1}$
نال ترکیب اے۔

\end{document}
